%==============================================================================
% Appendix Tables
%==============================================================================

%------------------------------------------------------------------------------
% Table A.1: Sample counts by treatment arm and outcome
%------------------------------------------------------------------------------
\begin{table}[htbp]
\centering
\caption{Sample counts by treatment arm and outcome}
\label{tab:counts}
\begin{tabular}{l ccccc}
\toprule
 & (1) & (2) & (3) & (4) & (5) \\
 & Personal SE belief & Trial belief & Vacc.\ intent & Link click & Vaccinated \\ \midrule
Control &   882 &   882 &   879 &   882 &   743 \\
Industry &   878 &   878 &   876 &   878 &   755 \\
Academic &   880 &   880 &   876 &   880 &   746 \\
Representative &   885 &   885 &   885 &   885 &   749 \\
All & 3,525 & 3,525 & 3,516 & 3,525 & 2,993 \\
\bottomrule
\end{tabular}
\tabnote{\linewidth}{Notes: Each cell reports the number of main-survey
respondents in the indicated treatment arm with non-missing values for the
indicated outcome. Column 3 (vaccination intention) has slightly fewer
observations due to missing responses to the intention question. Column 5
(vaccinated) reports follow-up survey respondents, reflecting attrition
between the main and follow-up surveys. See Table~\ref{tab:treat_effects}
for outcome definitions.}
\end{table}

\clearpage

%------------------------------------------------------------------------------
% Table A.2: Treatment effects on trust and relevance perceptions
%------------------------------------------------------------------------------
\begin{table}[htbp]
\centering
\caption{Treatment effects on trust and relevance perceptions}
\label{tab:trust_relevance}
\begin{tabular}{l cccc}
\toprule
 & (1) & (2) & (3) & (4) \\
 & Trust in trial & Relevance of trial & Trust in acad.\ study & Relevance of acad.\ study \\ \midrule
Industry            &      -0.449&       0.033&            &            \\
                    &     (0.102)&     (0.134)&            &            \\
Academic            &      -0.307&      -0.114&            &            \\
                    &     (0.100)&     (0.133)&            &            \\
Representative      &      -0.416&      -0.248&      -0.136&      -0.252\\
                    &     (0.101)&     (0.135)&     (0.103)&     (0.137)\\
Control/Academic mean&       5.751&       5.111&       5.957&       5.046\\
N                   &       3,520&       3,520&       1,762&       1,762\\
 \\
\bottomrule
\end{tabular}
\tabnote{\linewidth}{Notes: Each column is a separate OLS regression of the
indicated outcome on treatment arm indicators, with pre-specified controls
included but not shown. Robust standard errors in parentheses. Outcomes are
measured on a 0--10 scale. Columns 1--2 use all four arms with Control as the
omitted category. Columns 3--4 are restricted to Academic and Representative
arm respondents, as the academic study trust and relevance questions were only
asked of those arms; Academic is the omitted category so only the
Representative arm coefficient is reported. The base mean for columns 3--4 is
the Academic arm mean.}
\end{table}

\clearpage

%------------------------------------------------------------------------------
% Table A.3: Treatment effects on personal SE belief by prior flu vaccine experience
%------------------------------------------------------------------------------
\begin{table}[htbp]
\centering
\caption{Treatment effects on personal side-effect belief by prior flu vaccine experience}
\label{tab:hte_flu_vacc}
\begin{tabular}{l cccc}
\toprule
 & (1) & (2) & (3) & (4) \\
 & No vaccine & No reaction & Mild reaction & Severe reaction \\ \midrule
\input{../output/tables/het_flu_vacc_experience.tex} \\
\bottomrule
\end{tabular}
\tabnote{\linewidth}{Notes: Each column is a separate OLS regression of
personal side-effect belief ($\tilde{\Delta}$) on treatment arm indicators
(Industry, Academic, Representative), with Control as the omitted category.
The sample is split by prior flu vaccine experience: never vaccinated (column
1), vaccinated with no reaction (column 2), vaccinated with a mild reaction
(column 3), and vaccinated with a severe reaction (column 4). Pre-specified
controls are included but not shown. Robust standard errors in parentheses.
The bottom row indicates the subgroup for each column.}
\end{table}

\clearpage

%------------------------------------------------------------------------------
% Table A.4: Heterogeneous treatment effects by predicted trust/relevance index
%------------------------------------------------------------------------------
\begin{table}[htbp]
\centering
\caption{Heterogeneous treatment effects by predicted trust and relevance}
\label{tab:hte_pca}
\begin{tabular}{l cc}
\toprule
 & (1) & (2) \\
 & Personal SE belief & Vacc.\ intention \\ \midrule
Treatment: industry arm&      -3.300&      -0.001\\
                    &     (1.024)&     (0.012)\\
Treatment: academic arm&      -2.901&       0.031\\
                    &     (1.010)&     (0.012)\\
Treatment: personal arm&       0.369&       0.025\\
                    &     (1.021)&     (0.012)\\
Fitted values, penalized&      -1.311&       0.038\\
                    &     (1.961)&     (0.022)\\
Fitted values, penalized $\times$ Treatment: industry arm&      -2.985&       0.041\\
                    &     (1.779)&     (0.019)\\
Fitted values, penalized $\times$ Treatment: academic arm&      -1.982&       0.051\\
                    &     (1.820)&     (0.019)\\
Fitted values, penalized $\times$ Treatment: personal arm&      -2.439&       0.063\\
                    &     (1.885)&     (0.019)\\
Control mean        &      15.427&       0.069\\
N                   &       3,525&       3,516\\
 \\
\bottomrule
\end{tabular}
\tabnote{\linewidth}{Notes: Each column is a separate OLS regression of the
indicated outcome on treatment arm indicators (Industry, Academic,
Representative), a predicted trust/relevance index (\texttt{pca1\_hat}), and
interactions of each arm indicator with the index. \texttt{pca1\_hat} is the
first principal component of post-treatment trust in the clinical trial and
perceived relevance of the trial, predicted from pre-experiment covariates
using lasso selected in the control group. Rows labeled ``Fitted values,
penalized'' report the coefficient on \texttt{pca1\_hat}; interaction rows
report the coefficient on the arm-by-index interaction. Pre-specified controls
are included but not shown. Robust standard errors in parentheses. See
Table~\ref{tab:treat_effects} for outcome definitions.}
\end{table}
